\begin{tikzpicture}
\begin{axis}[spektrum,
    ymin=-0.15, ymax=1.2, ytick=\empty,
    ylabel={$|c_k|$},
    title={\textbf{Amplituden}}, title style={at={(0.85,1)}, anchor=south}]
  % Einhüllende
  \addplot[huelle, domain=0.6:10.5]   {sqrt(2)/(2*x)};
  \addplot[huelle, domain=-10.5:-0.6] {-sqrt(2)/(2*x)};
  \addplot[huelle, domain=0.85:10.5]  {1/x};
  \addplot[huelle, domain=-10.5:-0.85]{-1/x};
  \draw[blue, -latex, thin] (axis cs:4.2,0.62) node[right, font=\scriptsize]{$\frac{1}{\omega}$} to[bend right] (axis cs:1.65,0.61);
  \draw[blue, -latex, thin] (axis cs:5.6,0.36) node[right, font=\scriptsize]{$\frac{\sqrt{2}}{2\omega}$} to[bend right] (axis cs:3.1,0.23);
  % |c_k| = A_|k| / 2   (Ausnahme c_0 = A_0)
  \addplot[stab] coordinates {
    (-10,0.100) (-9,0.079) (-8,0) (-7,0.101) (-6,0.167) (-5,0.141)
    (-4,0) (-3,0.236) (-2,0.500) (-1,0.707)
    (0,0.800)
    (1,0.707) (2,0.500) (3,0.236) (4,0) (5,0.141)
    (6,0.167) (7,0.101) (8,0) (9,0.079) (10,0.100)};
\end{axis}
\end{tikzpicture}
\end{minipage}\hfill
\begin{minipage}[t]{0.49\linewidth}
\centering
\begin{tikzpicture}
\begin{axis}[spektrum,
    ymin=-2.8, ymax=2.8, ytick=\empty,
    ylabel={$\arg(c_k)$},
    xticklabel style={fill=white, inner sep=0.8pt},
    title={\textbf{Phasen}}, title style={at={(0.85,1)}, anchor=south}]
  % Hilfslinien bei Vielfachen von pi/4, Beschriftung rechts
  \pgfplotsinvokeforeach{-3,-2,-1,1,2,3}{
    \addplot[blue, dashed, thin, domain=-10.5:10.5, samples=2] {#1*pi/4};
  }
  \begin{scope}[every node/.style={blue, font=\tiny, anchor=west}]
    \node at (axis cs:10.5, 2.356) {$\frac{3\pi}{4}$};
    \node at (axis cs:10.5, 1.571) {$\frac{\pi}{2}$};
    \node at (axis cs:10.5, 0.785) {$\frac{\pi}{4}$};
    \node at (axis cs:10.5,-0.785) {$-\frac{\pi}{4}$};
    \node at (axis cs:10.5,-1.571) {$-\frac{\pi}{2}$};
    \node at (axis cs:10.5,-2.356) {$-\frac{3\pi}{4}$};
  \end{scope}
  % arg(c_k) = sgn(k) * phi_|k|
  \addplot[stab] coordinates {
    (-10,-1.571) (-9,2.356) (-8,0) (-7,0.785) (-6,-1.571) (-5,2.356)
    (-4,0) (-3,0.785) (-2,-1.571) (-1,2.356)
    (0,0)
    (1,-2.356) (2,1.571) (3,-0.785) (4,0) (5,-2.356)
    (6,1.571) (7,-0.785) (8,0) (9,-2.356) (10,1.571)};
\end{axis}
\end{tikzpicture}